\cotefile{157-1}
\scanpage{166}
\mspage{24}

Nous nous plaçons \inserted{\unsure{donc}} dans l'hypothèse où PS1 est satisfait (puisque b), c') et d) impliquent PS1), donc où \inserted{on est sous} les conditions équivalentes de 2.5. -- $\unsure{\mathscr X|\mathcal X}$ est loc. irréductible. Considérons \inserted{alors} la condition PS2, et prouvons qu'elle est impliquée tant par b), que par c'), et qu'elle implique \unsure{c'|a'}. Elle signifie que pour une double-flèche $a\overset{u}{\underset{v}{\unsure{=|\rightrightarrows}}}b$ dans $S$, \unsure{le|leur} \deleted{produit fibré} \inserted{\unsure{conoyau|noyau}} dans $S^\wedge$, ou ce qui revient au $=$ dans $\mathscr X^\wedge=S^\sim\subset S^\wedge$, est ``non vide''. On peut l'écrire sous la forme d'un produit fibré

\diagram{\begin{tikzcd}[ampersand replacement=\&]
a \arrow[r,"{(u,v)}"] \& b\times b \\
{\operatorname{Ker}(u,v)} \arrow[u] \arrow[r] \& b \arrow[u,"{\unsure{\mathrm{diag}|\mathrm{di}\,\Theta}}"']
\end{tikzcd}}

Ici par hyp. $a$ et $b$ sont des objets $0$-connexes de $\mathscr X^\wedge$ (comme ds $S$ l'étaient), \inserted{c) est vérifié ($\mathscr X$ loc. irréd. et ses \unsure{composantes} ttes $0$-connexes)} \deleted{Si $b\times b$ est $0$-connexe, alors si} donc $b\times b$ est \unsure{aussi} $0$-connexe, comme produit de deux objets \inserted{$0$-connexes} \unsure{sur une} \uline{$=$ \unsure{composante} connexe} ce qu'elle est \deleted{\unsure{est}} $0$-connexe\unsure{;} Comme $\mathscr X\times\mathscr X$ est loc. irréductible, il s'ensuit que le produit fibré est \deleted{\unsure{non}} vide (2.8.), donc $\mathrm{b)}\Longrightarrow\mathrm{d)}$. Prouvons $\unsure{\mathrm{d)}|\mathrm{c')}}\Longrightarrow\unsure{\mathrm{a)}|\mathrm{d)}}$ \unsure{i.e. si} \deleted{\unsure{les produits|que les produits} fibrés $a\times_{\unsure{\mathscr X|x}}b$ \unsure{($a,b$ ds $S$)} sont $0$-connexes, donc PS2 est satisfait,}

% archive sheet number: 166
